% Part: first-order-logic
% Chapter: syntax-and-semantics
% Section: substitution

\documentclass[../../../include/open-logic-section]{subfiles}

\begin{document}

\olfileid[pt]{fol}{syn}{sub}

\olsection{Substituição}

\begin{defn}[Substituição em um termo]
Definimos $\Subst{s}{t}{x}$, o resultado de \emph{substituir} $t$
por toda ocorrência de~$x$ em $s$, recursivamente:
\begin{enumerate}
\item \indcase{s}{c}{$\Subst{\indfrm}{t}{x}$ é simplesmente $s$.}

\item \indcase{s}{y}{$\Subst{\indfrm}{t}{x}$ é também simplesmente~$s$,
  desde que $y$ seja uma variável e $y \not\ident x$.}

\item \indcase{s}{x}{$\Subst{\indfrm}{t}{x}$ é~$t$.}

\item\indcase{s}{\Atom{f}{t_1, \dots, t_n}}{$\Subst{\indfrmp}{t}{x}$ é
  $\Atom{f}{\Subst{t_1}{t}{x}, \dots, \Subst{t_n}{t}{x}}$.}
\end{enumerate}
\end{defn}

\begin{defn}
Um termo~$t$ está \emph{!!{free for}} $x$ em $!A$ se nenhuma das
ocorrências livres de~$x$ em $!A$ ocorre no escopo de um quantificador que
liga uma variável em~$t$.
\end{defn}

\begin{ex} ~
\begin{enumerate}
\item $\Obj v_8$ está livre para $\Obj v_1$ em $\lexists[\Obj
  v_3]\Atom{\Obj A^2_4}{\Obj v_3,\Obj v_1}$

\item $\Obj f^2_1(\Obj v_1, \Obj v_2)$ \emph{não} está livre para $\Obj
  v_0$ em $\lforall[\Obj v_2]\Atom{\Obj A^2_4}{\Obj v_0,\Obj v_2}$
\end{enumerate}
\end{ex}

\begin{defn}[Substituição em !!a{formula}]
Se $!A$ é !!a{formula}, $x$~é !!a{variable}, e $t$~é um termo
!!{free for}~$x$ em~$!A$, então $\Subst{!A}{t}{x}$ é o resultado de
substituir $t$ por todas as ocorrências livres de~$x$ em~$!A$.
\begin{enumerate}
\tagitem{prvFalse}{\indcase{!A}{\lfalse}{$\Subst{\indfrm}{t}{x}$ é
    $\lfalse$.}}{}

\tagitem{prvTrue}{\indcase{!A}{\ltrue}{$\Subst{\indfrm}{t}{x}$ é
    $\ltrue$.}}{}

\item \indcase{!A}{\Atom{P}{t_1,\dots,
    t_n}}{$\Subst{\indfrm}{t}{x}$ é $\Atom{P}{\Subst{t_1}{t}{x},
    \dots, \Subst{t_n}{t}{x}}$.}

\item \indcase{!A}{\eq[t_1][t_2]}{$\Subst{\indfrmp}{t}{x}$ é
  $\Subst{t_1}{t}{x} = \Subst{t_2}{t}{x}$.}

\tagitem{prvNot}{\indcase{!A}{\lnot !B}{$\Subst{\indfrmp}{t}{x}$ é
    $\lnot \Subst{!B}{t}{x}$.}}{}

\tagitem{prvAnd}{\indcase{!A}{(!B \land
    !C)}{$\Subst{\indfrmp}{t}{x}$ é $(\Subst{!B}{t}{x} \land
    \Subst{!C}{t}{x})$.}}{}

\tagitem{prvOr}{\indcase{!A}{(!B \lor
    !C)}{$\Subst{\indfrmp}{t}{x}$ é $(\Subst{!B}{t}{x} \lor
    \Subst{!C}{t}{x})$.}}{}

\tagitem{prvIf}{\indcase{!A}{(!B \lif
    !C)}{$\Subst{\indfrmp}{t}{x}$ é $(\Subst{!B}{t}{x} \lif
    \Subst{!C}{t}{x})$.}}{}

\tagitem{prvIff}{\indcase{!A}{(!B \liff
    !C)}{$\Subst{\indfrmp}{t}{x}$ é $(\Subst{!B}{t}{x} \liff
    \Subst{!C}{t}{x})$.}}{}

\tagitem{prvAll}{
\indcase{!A}{\lforall[y][!B]}{$\Subst{\indfrmp}{t}{x}$
    é $\lforall[y][\Subst{!B}{t}{x}]$, desde que $y$ seja uma variável
    diferente de $x$; caso contrário $\Subst{\indfrmp}{t}{x}$
    é simplesmente $\indfrm$.}}{}

\tagitem{prvEx}{
\indcase{!A}{\lexists[y][!B]}{$\Subst{\indfrmp}{t}{x}$
    é $\lexists[y][\Subst{!B}{t}{x}]$, desde que $y$ seja uma variável
    diferente de $x$; caso contrário $\Subst{\indfrmp}{t}{x}$
    é simplesmente $\indfrm$.}}{}
\end{enumerate}
\end{defn}

\begin{explain}
Note que a substituição pode ser vácua: Se $x$ não ocorre de modo algum em $!A$,
então $\Subst{!A}{t}{x}$ é simplesmente~$!A$.

A restrição de que $t$ deve estar !!{free for}~$x$ em~$!A$ é necessária para
excluir casos como o seguinte. Se $!A \ident \lexists[y][x < y]$
e $t \ident y$, então $\Subst{!A}{t}{x}$ seria $\lexists[y][y <
  y]$. Neste caso a variável livre $y$ é ``capturada'' pelo
quantificador $\lexists[y]$ ao substituir, e isso é indesejável.
Por exemplo, gostaríamos que fosse o caso que, sempre que
$\lforall[x][!B]$ vale, $\Subst{!B}{t}{x}$ também vale. Mas considere
$\lforall[x][\lexists[y][x < y]]$ (aqui $!B$ é $\lexists[y][x <
  y]$). É uma sentença que é verdadeira a respeito, por exemplo, dos números naturais:
para todo número~$x$ há um número~$y$ maior que ele. Se
permitíssemos $y$ como uma possível substituição para~$x$, terminaríamos com
$\Subst{!B}{y}{x} \ident \lexists[y][y < y]$, que é falsa. 
Prevenimos isso exigindo que nenhuma das variáveis livres em~$t$
acabe sendo ligada por um quantificador em~$!A$.
\end{explain}

Frequentemente usamos a seguinte convenção para evitar notação incômoda: Se
$!A$ é !!a{formula} que pode conter a !!{variable}~$x$ livre, também
escrevemos~$!A(x)$ para indicar isso. Quando está claro qual $!A$
e~$x$ temos em mente, e $t$ é um termo (suposto estar livre para $x$
em $!A(x)$), então escrevemos $!A(t)$ como abreviação de $\Subst{!A}{t}{x}$. Assim,
por exemplo, poderíamos dizer, ``chamamos $!A(t)$ de uma instância
de~$\lforall[x][!A(x)]$.'' Com isso queremos dizer que se $!A$~é qualquer
!!{formula}, $x$~!!a{variable}, e $t$~um termo que está livre para~$x$
em~$!A$, então $\Subst{!A}{t}{x}$ é uma instância de~$\lforall[x][!A]$.

\end{document}
